\chapter*{Adnotare}

\par Negai Marin, \textit{``Erori vechi, protocol nou: analiza cauzelor eșecurilor de securitate din ecosistemul Model Context Protocol în condițiile adoptării accelerate''}, articol științific, Universitatea Tehnică a Moldovei, Facultatea Calculatoare, Informatică și Microelectronică, Chișinău, 2026. Articolul este alcătuit din introducere, patru capitole, concluzii, o bibliografie din 30 de surse și o anexă. Lucrarea conține 5 figuri și 5 tabele.

\vspace{0.3cm}
\par \noindent \textbf{Cuvinte-cheie:} Model Context Protocol, agenți AI, securitate software, injectare de prompt, analiza vulnerabilităților.
\vspace{0.3cm}

\par Model Context Protocol (MCP) a fost lansat la sfârșitul anului 2024 ca o modalitate comună de a conecta asistenții bazați pe modele de limbaj la fișiere, baze de date și servicii online. În decurs de un an, protocolul a fost preluat de majoritatea companiilor mari din domeniul AI, iar cercetătorii au publicat zeci de rapoarte despre servere vulnerabile și atacuri care ascund instrucțiuni în descrierile instrumentelor. Aceste probleme sunt prezentate de obicei ca un risc nou, specific AI. Articolul verifică dacă această interpretare este corectă sau dacă majoritatea eșecurilor sunt defecte software obișnuite, apărute pentru că ecosistemul a crescut mai repede decât a putut fi verificat.

\par Obiectivul lucrării a fost identificarea surselor eșecurilor de securitate din MCP. Au fost colectate \nRecords{} de vulnerabilități și incidente documentate public, dezvăluite între martie 2025 și februarie 2026, dintr-un catalog deschis de vulnerabilități și din publicațiile originale. Fiecare înregistrare a fost citită în sursa originală și clasificată pe două axe. Prima axă descrie natura defectului: clasic (ar exista și fără un model de limbaj), specific AI (funcționează doar pentru că modelul urmează un text în care nu ar trebui să aibă încredere) sau hibrid. A doua axă indică etapa în care a apărut defectul: proiectarea protocolului, implementarea, proiectarea integrării, lanțul de aprovizionare sau configurarea la instalare. Toate valorile au fost calculate de un script Python pe baza setului de date publicat.

\par Rezultatele arată că \pClassicComma\% dintre înregistrări sunt vulnerabilități clasice, precum injectarea de comenzi, traversarea căilor, falsificarea cererilor din partea serverului și lipsa autentificării, iar jumătate din totalul înregistrărilor sunt greșeli de programare în servere, clienți sau instrumente. Vulnerabilitățile specifice AI reprezintă \pNativeComma\% din eșantion și provin aproape în totalitate din proiectarea protocolului și a integrărilor. Niciuna dintre ele nu a primit un identificator CVE, deci instrumentele bazate pe baze de date de vulnerabilități nu pot observa această parte a riscului. Articolul propune un set de măsuri pentru dezvoltatori, operatori și instituții publice și compară rezultatele cu incidentul din octombrie 2026 care a afectat serviciile de guvernare electronică din Republica Moldova, unde o eroare de încredere asemănătoare a dus la expunerea datelor personale a aproximativ 450.000 de cetățeni.
